% 结构版本 v0.2（D015）；6.1–6.3 定稿于 V014（2026-09-26），6.4/6.5 仍为骨架。
% 状态：V014 定稿（用户对话逐句审定＋四轮网页 GPT 评审逐条仲裁后冻结，A013 预览—审定—写入模式）。仲裁要点：
% ①并发双事件 Cache 模型——查询在发射（P3:543）、插入在完成（P3:508-510），同键后完成者不重复插入不移位（insert_cache 已存在直接 return，P3:420-421）；𝒬(t) 队列状态与 𝒞(t) 驻留集合分离。
% ②B_DDR^CI 与官方 added 切割——added=partition+spill 在 Cache 模拟之前算出（本工作区亲验 P3:255-288），同方案两配置取值相同；命中改变 DDR 实际流量与竞争、不改逻辑统计量。
% ③250/60 均为池内公平共享，命中非必然更快；两来源收益（争用缓解/独立高带宽读池）＋关键链条件＋非单调结论。
% ④6.2.1 两类访问来源（结构性跨核重复＋spill 换入，三类 COPY_IN 同键空间）；"参与 Cache 查询"统一口径（超容量张量计入分母）。
% ⑤P3 单目标——对齐 1.5.3（V009"P3 不宣布额外搬运为优化目标"）与 1.3.3（V010 接受准则只看 Makespan）。
% ⑥三层区分（固定机制/描述状态/可控决策）唯一载体收至 6.3 章首，6.3 两小节化（原"缓存状态、固定机制与相关约束"并入 6.3.1，用户批准 2026-09-26，大纲已同步）。
% 第五章 V013（5.1–5.3）同期落稿后的记号对齐：mk_C 与 5.3 节 mk_B 同型（5.3 注释预留的第六章接口）；可行集合沿用 𝒟feas^B 不新造；𝒬(t)/𝒞(t) 规避第五章 C(t)（消费者集合）与第四章 Q_D（增强图）。
% 素材：题面 1.5/问题三/附录 B.4/D.5/D.6 逐字摘录（含"未找到"清单：命中固定延迟/初始状态/写路径分配/无 L2 基线产生方式题面未明说）＋官方评估器 P3 代码行号级核验＋2025A 国一第六章写法对照（融合式开篇、末问章轻公式重叙述、其第六章无机制示意图）＋2024A-26 精读接口警示（前问接口须逐式确认，落实于 6.3.1 映射链）。
% 边界（D039/D048）：6.1 为 L2 机制事实全章唯一展开处（2.4 仅定性引用、1.2.1 只留引入句）；6.2 只建访问关系/命中/带宽归属的形式化表达，规则引用 6.1 不复述。
% 延后项：F13 引用句随图定版补入（画法＝单次 COPY_IN 决策流：查询→Hit/Miss 两支→完成→若未驻留则插入→容量不足 FIFO 淘汰；COPY_OUT/spill-out 旁路；固定/可控/派生三框——05 备注）；同方案配对实验协议留 6.5.1；两配置搬运量统计口径留 6.5.4；"第五章 5.3 节"手写节号待第五章 label 体系稳定后复核。
% 符号首现（Needs-integration 待 3.2/TBL-003 收录；D/x/π_k/𝒟plan 沿用 4.2、mk 家族/𝒟feas^B 沿用 5.3 不重复登记）：a_j=(κ_j,β_j,s_j,f_j)、𝒬(t)、𝒞(t)、C_L2、h_j、B_DDR^CI、R_hit、mk_C、S^Cache_(k,i)、T^noL2/T^L2。术语候选：逻辑张量标识（题面"逻辑 tensor id"）、FIFO。
\section{共享 L2 Cache 下的联合优化模型}\label{sec:q3}

问题三将硬件场景在场景 B 的基础上进一步扩展：引入全核共享的只读 L2 缓存。决策对象、方案表示与合法性条件均沿用不变，变化的是方案到执行性能的映射——如 2.4 节所述，缓存行为是切图、核心分配与执行顺序共同塑造的访问流的后果，须作为调度决策的后果维度加以刻画。本章 6.1 节说明 L2 的机制事实，6.2 节建立缓存状态与性能映射模型，6.3 节给出 L2 条件下的调度优化模型，6.4 节说明求解框架的相应扩展，6.5 节报告无 L2 与只读 Cache 两种配置的对照实验。

\subsection{共享 L2 Cache 机制分析}

问题三在场景 B 的硬件结构上增加一片全核共享的只读 L2 缓存：容量 1048576 bytes（题面记为 1 MB，见表~\ref{tab:hw-config}），读带宽 250 bytes/cycle。L2 只服务读取，其读写路径划分与图~\ref{fig:hw-arch} 一致：COPY\_IN 类读取在发射时查询 L2（节点定义见 1.2.2 节）；而 COPY\_OUT 与 L1/UB 层面的缓存换出（spill）等写路径不经过 L2，仍直接访问 DDR。L2 读带宽与 DDR 总带宽是两个相互独立的带宽资源：命中的搬运进入 L2 读带宽池，未命中的搬运与全部其他核外搬运共享 DDR 总带宽，两池分别结算、互不占用；各池内部的并发搬运公平共享该池带宽，完成时刻随并发数变化，与 1.2.1 节 DDR 访问的结算语义一致。

L2 以逻辑张量标识（题面称逻辑 tensor id）为管理单位，按先进先出（First-In-First-Out，FIFO）方式组织。COPY\_IN 发射时以被读取张量的逻辑标识查询缓存：命中则直接从 L2 读取，不改变该张量在 FIFO 队列中的位置；未命中则照常从 DDR 读取，读取完成后触发插入检查，若该逻辑张量此时尚未驻留缓存且大小不超过容量，则将其整块插入，剩余容量不足时淘汰最早进入缓存的张量，直至容纳。由于查询键是逻辑张量标识，同一张量的全部 COPY\_IN 访问共享同一缓存条目，无论访问发生在哪个核心；一次未命中的读取完成并插入之后发射的后续访问，才可能因该条目的存在而命中。边界情形有两种：单张量大小超过 L2 容量时从不缓存，对它的访问恒为未命中、每次均从 DDR 读取；缓存从空开始，除容量触发的 FIFO 淘汰外不存在其他替换条件。

上述机制全部为题设给定的固定行为：容量、带宽与替换规则均不可配置，方案字段中也不存在任何缓存控制项（1.2.3 节）。
% TODO(F13, 图定版后补句)：单次 COPY_IN 的决策流与 FIFO 机制如图~\ref{fig:l2-fifo} 所示（Hit 支标注 FIFO 顺序不变；Miss 支含"DDR 读取→完成→若未驻留则插入→容量不足 FIFO 淘汰"；COPY_OUT/spill-out 旁路；固定/可控/派生三框）。素材就绪，registry 状态"计划"，画法备注随交付块转 05。

\subsection{L2 Cache 数据复用模型}

在上述机制基础上，本节从访问事件、缓存状态与带宽归属三个层面建立形式化表达，缓存规则引用 6.1 节、不再复述。

\subsubsection{共享输入数据与访问关系}

L2 的复用机会来自对同一逻辑张量的重复读取，本文主要关注两类来源。其一为结构性重复读取：场景 B 下同一 Task 内的子图边界不再插入 COPY\_IN（1.2.3 节），当核心分配将同一张量的多个消费方置于不同核心时，各核各自发起 COPY\_IN 读取该张量，这类访问由方案的跨核消费结构决定，属于可静态导出的结构信息，对应 1.3.2 节的多消费者张量特征。其二为缓存换入（spill 换入）：核内 L1/UB 容量不足时，评估程序在核内调度阶段自动插入换入型 COPY\_IN（1.4 节）。两类访问均以原张量的逻辑标识为键，共同构成 L2 的服务对象。各次访问发生的时刻则由执行时序决定，只能在评估过程中获得。

将多核执行中参与 L2 Cache 查询的全部 COPY\_IN 访问记为事件集合，每次访问 $a_j=(\kappa_j,\beta_j,s_j,f_j)$：$\kappa_j$ 为被读取张量的逻辑张量标识，$\beta_j$ 为该次访问的字节数，$s_j$ 为发射时刻，$f_j$ 为执行模拟中内生确定的完成时刻。同一键的多次访问构成潜在的复用机会；一次访问是否真正命中，不仅取决于此前同键访问是否已完成并驻留缓存，还取决于期间是否已被淘汰，这是下一小节刻画的内容。

\subsubsection{Cache 命中与数据搬运}

缓存状态随执行过程演化：记 $\mathcal{Q}(t)$ 为时刻 $t$ 的 FIFO 缓存队列状态，包含各驻留逻辑张量、条目字节数及其插入顺序；$\mathcal{C}(t)$ 为其对应的驻留张量集合，评估从空缓存 $\mathcal{Q}(0)=\varnothing$ 开始。第 $j$ 次访问的命中判定为
\[ h_j=\begin{cases}1, & \kappa_j\in\mathcal{C}(s_j^{-}),\\ 0, & \text{其他},\end{cases} \]
其中 $s_j^{-}$ 表示发射时刻前一瞬时。对 $\beta_j$ 超过 L2 容量 $C_{\mathrm{L2}}$ 的张量，官方机制从不将其插入缓存，因此其对应访问恒有 $h_j=0$。

缓存命中判定与队列状态演化分别发生在发射事件与完成事件：未命中的访问在其完成时刻 $f_j$ 触发插入检查，若该逻辑张量此时尚未驻留缓存且大小不超过容量，则将其整块插入，容量不足时按 FIFO 淘汰队头条目；若同键条目已由其他先完成的访问插入，则不重复插入，也不改变其队列位置；命中同样不改变队列顺序。$\mathcal{Q}(t)$ 的精确更新按官方评估程序的离散事件语义计算，本文不作静态近似（与 1.3.3 节的信息分层一致）。

命中与否决定该次访问的数据来源与带宽归属：命中访问进入 L2 读带宽池结算，未命中访问从 DDR 读取 $\beta_j$ 字节、与全部其他核外搬运共同竞争 DDR 总带宽。据此定义由 DDR 实际服务的访问字节量
\[ B_{\mathrm{DDR}}^{\mathrm{CI}}=\sum_{j}(1-h_j)\,\beta_j, \]
它刻画参与 Cache 查询的 COPY\_IN 中仍需由 DDR 实际服务的数据规模，与命中率共同描述 L2 对读取流量的分流效果。须将其与官方指标“总额外数据搬运量”区分：后者由方案诱导的跨 Task 边界搬运与 L1/UB 层面的缓存换入换出决定，在缓存模拟之前由方案的静态结构算出，不因命中而扣减——对同一方案，无 L2 与只读 Cache 两种配置下该指标取值相同。命中所改变的是 DDR 实际服务的流量与带宽竞争，而非这一逻辑统计量；两配置搬运量对照的具体口径见 6.5.4 节。在访问集合上统计字节命中率
\[ R_{\mathrm{hit}}=\frac{\sum_{j} h_j\,\beta_j}{\sum_{j}\beta_j}, \]
即 1.4 节 Cache 命中率指标的形式化表达，分母为参与 Cache 查询的全部 COPY\_IN 访问字节（含缓存换入重载）。

\subsubsection{L2 对 DDR 访问需求及带宽竞争的影响}

L2 对 DDR 访问需求的作用是直接的：每次命中将 $\beta_j$ 字节的搬运需求从全核共享的 DDR 带宽池转移至独立的 L2 读带宽池，池内其余并发搬运可分得更多带宽、更早完成，从而缓解带宽争用。此外，L2 读带宽上限（250 bytes/cycle）高于 DDR 总带宽（60 bytes/cycle）且两池独立，命中访问通常可获得更高的数据服务能力；但两池内部均为公平共享，单次搬运的实际完成时间仍取决于该时刻池内的并发搬运数量，不能断言命中访问必然更快完成。

命中率与总体任务执行时间之间不存在必然的单调关系。命中的时间收益有上述两个来源，其成立条件各不相同：缓解 DDR 争用的收益只有在相关时段存在显著争用时才明显；经独立读池获得更高服务能力的收益则受 L2 池并发程度限制。二者最终能否缩短 Makespan，还取决于相关搬运能否影响最终的任务完成时刻；而调度决策的变化同时改变计算并行结构与访问时刻，为提高命中率而作的聚合式改动可能以牺牲并行度与负载均衡为代价。因此，本文不以静态命中率近似裁决候选方案的优劣，而以含 L2 的官方评估结果为准（见 6.3.2 节）；缓存收益的量化采用同一方案的两配置对照，实验协议见 6.5.1 节。

\subsection{L2 Cache 条件下的调度优化模型}

给定方案，L2 的机制参数固定，缓存状态由官方执行语义唯一派生：问题三没有新增任何方案层面的决策变量或约束，新增的是从方案到执行性能的缓存状态映射。本节先说明相对问题二模型的这一扩展（6.3.1 节），再给出目标函数与 L2 感知的评价方式（6.3.2 节）。

\subsubsection{对问题二模型的扩展}

问题三的可行方案集合与问题二完全一致：决策仍为 $D=(x,\{\pi_k\})$（记号沿用 4.2.1 节，两方案字段见 1.2.3 节），第五章 5.3 节的合法集合与执行可行性判定全部沿用，L2 不新增方案合法性约束——缓存机制只改变执行语义与性能，不改变合法方案的集合。扩展发生在决策到性能的映射链上：切图与核心分配首先确定跨核消费结构及由此产生的结构性读取；核内调度在容量不足时进一步产生缓存换入，形成完整的 COPY\_IN 访问集合；执行顺序与带宽争用共同确定各次访问的发射与完成时刻。访问事件经 6.2.2 节的命中与更新规则确定缓存动态，经带宽归属确定执行时间线，最终决定总体任务执行时间。求解算法可利用的缓存相关信息（如多消费者张量结构、重复读取的字节规模）只是服务于搜索的结构性先验，其利用方式属于算法设计，见 6.4 节。

由此，L2 容量与 L1/UB 容量的角色不同：L1/UB 容量直接参与核内执行可行性判断及缓存换入换出的生成，并由此影响方案的执行语义与性能；L2 容量则不构成对方案的约束条件，其容量边界完全由固定的插入与 FIFO 淘汰机制处理，经命中率与执行时间线作用于方案的性能评价。

\subsubsection{目标函数与 L2 感知评价}

问题三的优化目标仍为总体任务执行时间。记 $\mathrm{mk}_C(D)$ 为决策 $D$ 对应方案在问题三官方执行语义（含 L2）下评估所得的总体任务执行时间（场景标记与 5.3 节 $\mathrm{mk}_B$ 同型），问题三的优化问题在问题二的可行方案集合 $\mathcal{D}_{\mathrm{feas}}^{B}$（记号沿用 5.3 节，其构成见 6.3.1 节）上定义：
\[ \min_{D\in\mathcal{D}_{\mathrm{feas}}^{B}}\ \mathrm{mk}_C(D). \]
总额外数据搬运量与 Cache 命中率作为辅助报告与机理分析指标，不进入搜索的接受准则（1.3.3 节）。L2 的引入不需要也不应当修改目标函数的形式：缓存收益经由官方评估的执行模拟直接进入总体任务执行时间，以命中率替代或加权改写时间目标反而偏离题设的评价语义。

L2 通过两种方式进入评价。其一作为评价环境：候选方案的生成与比较在含 L2 的官方评估语义下进行，命中判定与搬运归属由评估过程确定，求解算法在该环境中利用缓存相关结构信息引导搜索，具体见 6.4 节。其二作为对照对象：题面要求以无 L2 与只读 Cache 两种配置的对照衡量缓存影响；为隔离缓存机制本身的作用，本文采用同一方案的配对评价，定义相对无 L2 基线的加速比——对测试用例 $i$ 的 $k$ 核配置，
\[ S^{\mathrm{Cache}}_{k,i}(D)=\frac{T^{\mathrm{noL2}}_{k,i}(D)}{T^{\mathrm{L2}}_{k,i}(D)}, \]
其中 $T^{\mathrm{noL2}}_{k,i}(D)$ 与 $T^{\mathrm{L2}}_{k,i}(D)$ 分别为方案 $D$ 在测试用例 $i$、$k$ 核配置下于无 L2 与只读 Cache 两种配置中评估所得的 Makespan（$T$ 沿用 1.4 节实验统计记号；两配置分别对应 $\mathrm{mk}_B$ 与 $\mathrm{mk}_C$ 在该实例与核数下的取值）；测试集上的平均与 1.4 节口径一致，取逐用例比值的算术平均。该指标是同一方案在两种硬件配置下的配对评价指标，不是优化目标；配对实验协议见 6.5.1 节。

\subsection{L2 感知调度优化算法}
\pending{内容 TODO：依据题面、最终采用的实现及适用证据补充本节；具体算法、公式和结果尚未填写。}

\subsubsection{算法思想与选择依据}
\pending{内容 TODO：依据题面、最终采用的实现及适用证据补充本节；具体算法、公式和结果尚未填写。}

\subsubsection{调度策略及与问题二的关系}
\pending{内容 TODO：依据题面、最终采用的实现及适用证据补充本节；具体算法、公式和结果尚未填写。}

\subsubsection{算法流程、伪代码与停止条件}
\pending{内容 TODO：依据题面、最终采用的实现及适用证据补充本节；具体算法、公式和结果尚未填写。}

\subsubsection{算法复杂度与计算预算}
\pending{具体策略 TODO，以最终实现为准。}

\subsection{实验结果与分析}
\pending{内容 TODO：依据题面、最终采用的实现及适用证据补充本节；具体算法、公式和结果尚未填写。}

\subsubsection{实验设置与无 L2 / 有 L2 配对协议}
\pending{固定用例、核数和必要的方案标识，区分机制对照与算法对照。}

\subsubsection{1～5 核下无 L2 与只读 L2 的平均加速比}
\pending{两组曲线均明确相对单核基准的定义。}

\subsubsection{无 L2 与只读 L2 的 Makespan 对比及相对加速比}
\pending{明确该相对加速比与 6.5.2 中指标的区别。}

\subsubsection{Cache 字节命中率与 DDR 数据搬运量}
\pending{区分总搬运量、总额外搬运量及其统计口径。}

\subsubsection{L2 感知策略对比与消融}
\pending{区分硬件缓存本身的收益和算法利用缓存的额外效果。}

\subsubsection{不同计算图的收益差异与机理分析}
\pending{结合访问、搬运、等待或调度记录解释受益及不受益的案例；没有相应证据时保留待验证解释。}

\subsubsection{求解成本与算法结果小结}
\pending{内容 TODO：依据题面、最终采用的实现及适用证据补充本节；具体算法、公式和结果尚未填写。}
